\section{Introduction}\label{sec:introduction}

For a finite set \(X\subset\R^2\), write
\[
 u(X)=\#\bigl\{\{x,y\}\subset X:\ \norm{x-y}_2=1\bigr\},
 \qquad
 u(n)=\max_{\substack{X\subset\R^2\\ |X|=n}}u(X).
\]
Thus pairs are unordered and their endpoints are distinct. We prove a fixed
power saving over the exponent \(4/3\).

\begin{theorem}\label{thm:main}
There are absolute constants \(0<C<\infty\) and \(1\le\beta<4/3\) such that
\(u(n)\le Cn^\beta\) for every nonnegative integer \(n\).
\end{theorem}

\Cref{thm:main} gives a positive answer to the question whether this upper
exponent can be lowered by a fixed amount: equivalently,
\(u(n)=O(n^{4/3-\delta})\) for an absolute \(\delta>0\).
All constants are independent of the configuration.
